% Integration fragment: no documentclass.
% Insert before Lemma lem:semiprime55 (or before the first HB interface).
% Uses existing bibliography key Lichtman2023 and ordinary lemma/proof.
% No distribution estimate for primes to large moduli is used in this lemma.

\subsection{Siegel--Walfisz stability of the individual rough factors}
\label{sec:rough-sw-stability}

The fact that an individual Heath--Brown variable has coefficient
$1$ or $\mu$ does not by itself check the additional roughness
restriction. The following elementary stability argument supplies
that check, including the auxiliary coprimality and the Mellin twists.

\begin{lemma}[Rough Siegel--Walfisz stability]
\label{lem:rough-sw}
Fix $\xi>0$ and $B\ge0$. Let
\[
 x^\xi\le U\le x,\qquad z_0=x^{1/(\log\log x)^3},
\]
and let $I\subset[U,2U]$ be any interval.
Let $e(n)$ be either $1$ or $\mu(n)$.
Suppose $h$ is a smooth function on $[U,2U]$ such that
\[
 |h(v)|\le1,\qquad
 \int_U^{2U}|h'(v)|\,dv\ll(\log x)^B,
 \qquad |t|\le(\log x)^B.
\]
Set
\[
 \alpha(n)=e(n)\mathbf1_{P^-(n)>z_0}\mathbf1_I(n)
                  h(n)n^{it}.
\]
For every $A>0$, all integers $D,q\ge1$, and every residue $r$
with $(r,q)=1$, one has
\begin{equation}\label{eq:rough-sw-full}
 \left|
 \sum_{\substack{n\sim U\\n\equiv r\pmod q\\(n,D)=1}}\alpha(n)
 -\frac1{\varphi(q)}
       \sum_{\substack{n\sim U\\(n,Dq)=1}}\alpha(n)
 \right|
 \ll_{A,B,\xi}\frac{U\,\tau(D)}{(\log x)^A}.
\end{equation}
The implied constant is uniform in $I,D,q,r,t$ and in the stated
bounds for $h$. Consequently these sequences meet
\cite[Definition 2.1 and Proposition 5.2]{Lichtman2023}'s
Siegel--Walfisz hypothesis. Fixed powers of $\log n$ are also
allowed after normalization by the corresponding power of
$\log(2U)$.
\end{lemma}

\begin{proof}
We first prove the assertion for $D=1$, $h=1$, and $t=0$,
uniformly for every interval $I\subset[U,2U]$.
The classical inputs are the ordinary Siegel--Walfisz estimate
for the unrough sequence $\mu$ and the fundamental lemma of
the one-dimensional sieve. The needed auxiliary-coprimality
version of the former can itself be reduced to the ordinary
version as follows.
For any integer $E\ge1$ the exact convolution identity is
\[
 \mu(n)\mathbf1_{(n,E)=1}
       =\sum_{\substack{v\mid n\\v\mid E^\infty}}\mu(n/v).
\]
On a scale $V$, split this at $v=V^{1/4}$.
For the small $v$, apply ordinary Siegel--Walfisz to $n/v$;
the sum of the factors $1/v$ is at most
$\prod_{p\mid E}(1-1/p)^{-1}\le\tau(E)$.
For the tail use
\[
 \sum_{\substack{v>V^{1/4}\\v\mid E^\infty}}\frac1v
 \le V^{-1/8}\prod_{p\mid E}(1-p^{-1/2})^{-1}
 \le V^{-1/8}\tau(E)^2.
\]
Both its progression contribution and its normalized coprime
mean are therefore $O(V^{7/8}\tau(E)^2)$, by bounding them
with the total absolute mass. Terms $(v,q)>1$ contribute zero
in the primitive residue and its coprime mean.
Thus the auxiliary-coprimality estimate has the fixed divisor
factor $\tau(E)^2$, uniformly on subintervals of $[V,2V]$
and in each fixed logarithmic range of $q$.
This uses no rough Siegel--Walfisz assertion.

\paragraph{Small moduli and the fundamental-lemma error.}
Fix a constant $B_0$ and first take $q\le(\log x)^{B_0}$.
Put $L=U^{1/4}$. For large $x$ one has $z_0<L$, and
\[
 s_0=\frac{\log L}{\log z_0}
       \ge\frac{\xi}{4}(\log\log x)^3.
\]
Take upper and lower fundamental-lemma weights $\lambda^\pm_\ell$
of level $L$, supported on squarefree divisors of $P(z_0)$,
with $|\lambda^\pm_\ell|\le1$, and write
\[
 \rho^\pm(n)=\sum_{\ell\mid n}\lambda^\pm_\ell.
\]
They satisfy
$\rho^-(n)\le\mathbf1_{P^-(n)>z_0}\le\rho^+(n)$.
Endpoint conventions at $z_0$ can be made identical throughout.

In a primitive residue class $r\bmod q$ the local density at a
sifting prime is $1/p$ if $p\nmid q$ and zero otherwise.
Deleting the primes dividing $q$ preserves the uniform
one-sided dimension-one hypothesis of the fundamental lemma.
The main-term difference of the two weights is therefore
$O(V_q(z_0)e^{-s_0})$, with $V_q(z_0)\le1$.
Counting multiples of $\ell$ in $I\cap\{n\equiv r\pmod q\}$
has an error $O(1)$ for each $(\ell,q)=1$, and none occur
when $(\ell,q)>1$. Hence
\begin{equation}\label{eq:rough-beta-gap}
 \sum_{\substack{n\in I\\n\equiv r\pmod q}}
              (\rho^+(n)-\rho^-(n))
 \ll \frac Uq e^{-s_0}+L
 \ll_C U(\log x)^{-C}
\end{equation}
for every fixed $C>0$.
The same bound holds for the average over the primitive residue
classes. This controls the normalized coprime main sum as well.
It is uniform in the endpoints of $I$ and in $q$ in the indicated
range.

Since $|e(n)|\le1$,
\[
 \big|e(n)\mathbf1_{P^-(n)>z_0}-e(n)\rho^+(n)\big|
       \le\rho^+(n)-\rho^-(n).
\]
Thus \eqref{eq:rough-beta-gap} justifies replacing the rough
indicator by $\rho^+$ in both terms of the desired discrepancy.
This absolute-error argument is necessary when $e=\mu$;
no positivity of $\mu$ is asserted.

\paragraph{The ordinary Möbius Siegel--Walfisz main sum.}
For $e=\mu$, the replacement progression sum becomes
\[
 \sum_{\substack{\ell\le L\\(\ell,q)=1}}
   \lambda^+_\ell\mu(\ell)
 \sum_{\substack{m\in I/\ell\\
             m\equiv r\bar\ell\pmod q\\(m,\ell)=1}}\mu(m),
\]
using the exact identity
$\mu(\ell m)=\mu(\ell)\mu(m)\mathbf1_{(m,\ell)=1}$
for squarefree $\ell$.
The normalized primitive-class average is the identical outer
sum with the inner sum replaced by
\[
 \frac1{\varphi(q)}
       \sum_{\substack{m\in I/\ell\\(m,\ell q)=1}}\mu(m).
\]
Here $U/\ell\ge U^{3/4}\ge x^{3\xi/4}$, so
$\log(U/\ell)$ is comparable with $\log x$ and
$q\le(\log x)^{B_0}$ lies in a fixed logarithmic
Siegel--Walfisz range for every inner sum.

Use the unrough auxiliary-coprimality estimate just derived:
with the fixed divisor exponent $C_0=2$ and any requested
logarithmic saving $C$ its inner discrepancy is
\[
 \ll_{C,B_0,\xi}
       \frac{(U/\ell)\tau(\ell)^{C_0}}{(\log x)^C}.
\]
The same estimate holds on arbitrary subintervals of
$[U/\ell,2U/\ell]$ by subtracting endpoint estimates.
Summing over $\ell$ and using the fixed divisor moment
\[
 \sum_{\ell\le L}\frac{\tau(\ell)^{C_0}}{\ell}
                   \ll(\log L)^{C_1}
\]
gives $O_A(U(\log x)^{-A})$ after choosing
$C>A+C_1+1$. The divisor exponent here is the fixed exponent
in the classical auxiliary-coprimality input; it is not
allowed to grow with $\ell$.
This calculation explains explicitly the auxiliary condition
$(m,\ell)=1$ introduced by the identity for $\mu(\ell m)$.

For $e=1$, expand $\rho^+$ in the same way without $\mu$.
For each $(\ell,q)=1$, counts of $m\in I/\ell$ in any two
residue classes modulo $q$ differ by $O(1)$.
Thus their difference from the primitive-class average
is also $O(1)$, and summing the weights costs $O(L)$.
Together with \eqref{eq:rough-beta-gap}, this proves the
unweighted assertion for all small $q$ in both cases.

\paragraph{Large moduli.}
For every sequence under consideration with absolute value
at most one, the discrepancy without auxiliary $D$ is trivially
\[
 \ll \frac Uq+1+\frac U{\varphi(q)}
       \ll Uq^{-1/2}+1.
\]
Choose $B_0>2A+4$. For $q>(\log x)^{B_0}$ this is
$O_A(U(\log x)^{-A})$, since $U\ge x^\xi$.
Together with the small-modulus proof, this gives the estimate
for all $q\ge1$, uniformly on every interval $I$.

\paragraph{An arbitrary auxiliary integer.}
Return to any $D\ge1$.
Primes at most $z_0$ dividing $D$ cannot divide a rough $n$.
The number of remaining integers of $[U,2U]$ removed by
$(n,D)=1$ is at most
\[
 \sum_{\substack{\ell\mid D\\\ell>z_0}}
           \left(\frac U\ell+1\right)
       \le\left(\frac U{z_0}+1\right)\omega(D).
\]
The change in the progression term plus its normalized
coprime mean is at most twice this bound, because the
coefficients are bounded by one and $1/\varphi(q)\le1$.
For every fixed $A$,
\[
 \left(\frac U{z_0}+1\right)\omega(D)
        \ll_{A,\xi}\frac{U\,\tau(D)}{(\log x)^A}.
\]
This is uniform for arbitrarily large $D$, and proves the
complete auxiliary-coprimality assertion. It does not
impose a hidden size restriction on that auxiliary integer.

\paragraph{Mellin twists and smooth coefficients.}
The estimates just obtained hold for every subinterval
of $[U,2U]$.
On this interval the function $h(v)v^{it}$ has total variation
\[
 \ll \int_U^{2U}|h'(v)|\,dv
                 +|t|\int_U^{2U}\frac{dv}{v}
 \ll(\log x)^B.
\]
Partial summation therefore costs at most a fixed power of
$\log x$. Apply the untwisted estimate with a correspondingly
larger arbitrary saving to obtain \eqref{eq:rough-sw-full}.
The interval indicator has only two endpoint jumps.
A normalized fixed power of $\log v$ has bounded variation
here, so the final assertion follows in the same manner.
\end{proof}

\paragraph{Use in the Heath--Brown arguments.}
In the preceding or following factorization arguments, take
$\xi=2\zeta$ for a selected individual variable with
$x^{2\zeta}\le U\le x^{1/3+\zeta/2}$.
Its coefficient after Mellin separation is precisely one of
the sequences in Lemma~\ref{lem:rough-sw}, after harmless
normalization of logarithms.
One must select an individual $1$ or $\mu$ slot; this lemma
does not assert Siegel--Walfisz for an arbitrary grouped
convolution.

Before Mellin separation, discard dyadic Heath--Brown boxes
which do not meet the original prime-variable interval.
If an original factor lies in $[R,2R]$, the product of its
at most six dyadic lower endpoints is then at least
$R/2^6$ and at most $2R$.
These group-scale inequalities remain true for the retained
boxes after Mellin separation.
They justify the assertion that, in the absence of medium
slots, each original group of size at least $N^a/2$
must contain a long slot when $12\zeta<a/2$.
At most two long slots can occur globally because three
would have product exceeding $x^{1+3\zeta/2}$.
